\odiseachapter{laertes}{Laertes}\label{cap:laertes}

El último canto de la \emph{Odisea} es un pastiche que incluye tres pasajes diferentes.

El primero es una segunda visita al Hades, con la excusa de seguir a las almas de los pretendientes, conducidas en tropel por Hermes al prado de los asfódelos. El poema sigue insistiendo en la imagen del rebaño sacrificado, que incluso después de muertos precisa de un pastor para guiarlos al averno. No es que Homero se interese mucho por la ex-alegre pandilla de gorrones. Toda la excursión al ultramundo es una excusa para presentarnos una conversación entre Agamenón y Aquiles, en la que el primero sigue quejándose ---claramente le sentó fatal que lo asesinaran--- y de paso describe, no con poca envidia, las honras fúnebres dedicadas al segundo.

Tampoco, por supuesto, se nos cuenta nada más de las esclavas, vilmente asesinadas. ¿Es ese ninguneo una prueba de «machismo»? Al margen de las dificultades de proyectar la ética y el lenguaje moderno a una civilización separada de la nuestra por tres milenios, me parece que no es el caso. Si excluimos a Odiseo, que, todo hay que decirlo, tiene algo de \emph{Übermensch}, me atrevo a decir, y así he tratado de reflejarlo en este libro, que las auténticas protagonistas de la \emph{Odisea} son mujeres. Algunas divinas ---Atenea, Circe, Calipso---, otras humanas ---Penélope, Nausícaa, Arete, Clitemnestra--- y alguna, como Helena, en la liga especial de las semidiosas. Pero ninguna de ellas es un florero, ninguna se pliega a los deseos masculinos, ninguna se limita a actuar como comparsa, todas se caracterizan por su firme voluntad y en la mayoría de los casos ---Penélope, Nausícaa, Arete, incluso las divinas Circe y Calipso--- por una calidad moral que ya querrían tener los hombres. En la división olímpica, aunque Zeus hace algún cameo, quien lleva absolutamente la voz cantante es Atenea. El lector simpatiza incluso con Clitemnestra, sobre todo en la versión en la que su venganza está motivada por el asesinato de su hija y resulta difícil no acabar suspirando por Helena, como aquellos bárbaros de hace tres milenios.

Así que el trato que se les da a las muchachas se entiende, precisamente, en la clave de que son esclavas. Por eso las ahorca Telémaco, por eso nada se nos dice de su viaje al ultramundo. Como tales esclavas no tienen derecho, ya no digamos a un indulto más que justo dada la levedad ---por no decir inexistencia--- de sus pecados, sino ni siquiera a una muerte «digna» ---al menos a los ojos de los que las ejecutan---. Y en ese mismo espíritu, tampoco tienen derecho a que el Poeta se fije en sus sombras.

Son esclavas, anónimas, despreciadas, ignoradas por sus amos. Dicho sea de paso, el cabrero Melantio, que, en última instancia, es otro siervo, tampoco asoma por ninguna parte, aunque solo sea para protestar por su nariz, sus orejas y sus verijas extirpadas.

Afortunadamente, en estos tiempos modernos no hay esclavitud, ni todo un continente ---África--- donde la renta per cápita anual, el acceso a la energía, al agua potable y a servicios básicos de enseñanza y salud son tan inferiores a los del primer mundo que la fortuna del gran magnate que todos conocemos rivaliza con el producto anual de las mayores economías del continente. Tampoco hay guerras innecesarias y crueles, ni genocidios, ni desigualdades entre los adinerados y los desposeídos, ni millones de personas dispuestas a jugarse la vida por el espejismo de una vida mejor, es decir, esclavos que se suben a la patera como las muchachas de la casa de Ulises se subían al lecho de los pretendientes, porque no tienen nada mejor a lo que agarrarse.

Quizás por eso nos escandalizamos tanto frente al abuso. Tanto, que ninguna película lo muestra, porque barrer la ignominia bajo la alfombra es algo a lo que estamos muy bien acostumbrados.

El segundo fragmento del canto XXIV es la visita a Laertes. Odiseo no las tiene todas consigo cuando va a ver a su padre. La ausencia ha sido demasiado larga y tiene miedo de que no le reconozca. Les dice a sus criados:

\begin{versocita}
«Ahora marchaos y entrad en la bien puesta casa vosotros,\\
y para almuerzo matáis enseguida al marrano más gordo;\\
que yo iré a ver si mi padre, cuando me tenga a sus ojos,\\
me reconoce y se da cuenta aún de quien soy, o del todo\\
me desconoce, que muy largo tiempo tomó mi retorno».
\end{versocita}

Ulises se acerca hasta la casita donde vive su padre ---nunca se nos explica por qué Laertes no reside en palacio y ha optado por una humilde morada en el campo--- y lo encuentra anciano y desaliñado:

\begin{versocita}
Halló, pues, solo a su padre en aquel belguarnido viñedo\\
escardillando una planta; vestía un jubón sucio y viejo\\
sin apostura y zurcido, y llevaba polainas de cuero,\\
cosidas, como en torzal, para no rasguñarse los remos,\\
y guante tosco en las manos para las zarzas, y un yelmo\\
de piel de cabra en su testa; así incubaba su duelo.
\end{versocita}

Como es habitual en Homero, la atención al detalle no tiene nada que envidiar a los mejores narradores modernos. Y por una vez, el milerrante, biensagaz, milardido Odiseo se queda sin trucos, abrumado por la pena:

\begin{versocita}
[...] detrás de un alto peral dejó ir una lágrima y quieto\\
allí se estuvo dudando en el alma y en su pensamiento\\
si ir a besar y a abrazar a su padre y contarle al momento\\
cómo se obró su regreso y cómo llegó al patrio suelo,\\
o preguntar él primero y probarlo en hablas y verbos.
\end{versocita}

Pero Ulises no puede evitar seguir con sus tretas, incluso cuando se trata de su padre:

\begin{versocita}
Tal cavilando, en la duda, pensó que sería más bueno\\
entrarle en voz zahiriente para probarlo primero.\\
Esto pensando, derecho fue a él el divino Odiseo.
\end{versocita}

Así que se va hacia Laertes y le espeta:

\begin{versocita}
Viejo, saber no te falta a ti para dar un buen cuido\\
a un huerto de estos, que todo lo tienes muy bien atendido,\\
y no se ve planta alguna, ni higuera, ni vid, ni aun olivo,\\
y ni un peral, ni un arriate por este vergel en descuido.\\
Mas te diré cosa más, y que a enojo no caiga en tu espíritu:\\
cuidado bueno no tienes en atenderte a ti mismo,\\
que, a más de en triste vejez, vas sucio y muy mal vestido.
\end{versocita}

Después de imprecarlo así, continúa, haciéndose el longuis:

\begin{versocita}
Y aún otra cosa —y verdad— porque quiero saberla te pido:\\
dime si es cierto que estamos en Ítaca, que eso me dijo\\
uno que ha nada encontré, al venir yo hacia aquí, en el camino,\\
y andaba un poco aturdido y decirme al final no ha sabido\\
—que ni quería escucharme— cuando por un huésped mío\\
le pregunté, que no sé si por caso él aún está vivo\\
o si murió y su morada la tiene con Hades vecino.
\end{versocita}

El muymentiroso asegura entonces que el huésped en cuestión «dijo más de una vez que nació de Laertes, el hijo de Arcisio», y el pobre viejo responde, llorando:

\begin{versocita}
Sí, extranjero, es así, en la tierra tú estás que buscabas,\\
la misma tierra que hoy tienen soberbios hombres sin tasa;\\
vanos, por más que mil dieras, serán tus regalos de gracia.\\
Si en este pueblo de Ítaca vivo a tu huésped hallaras,\\
trocando doñas debidas, muy bien regalado marcharas\\
y en buena hospitalidad, como manda la ley que a ella ata.
\end{versocita}

Como era de esperar, a Laertes le falta tiempo para pedirle al extranjero detalles sobre el paradero del supuesto huésped ---es decir el hijo que le falta desde hace dos décadas--- y también le pregunta por su identidad. Ulises bien podía haber dado por terminado el jueguito, pero sigue mintiendo, el muy bellaco. Se inventa una ---otra más--- falsa identidad y le da largas sobre el paradero de Odiseo, provocando la desesperación de Laertes.

\begin{versocita}
Dijo; y cayó sobre el viejo oscura una nube de duelo;\\
y, en ambas manos cogiendo del polvo trasceniciento,\\
sobre su cana cabeza lo echó en apretado lamento.
\end{versocita}

Por fin, Ulises parece ablandarse, se identifica, abraza a su padre. Ah, pero se demora poco tiempo en sentimentalismos y enseguida pasa a los negocios, informando a su progenitor de la matanza de los pretendientes.

\begin{versocita}
Y ya a Odiseo le hirvió el corazón, y un acre veneno\\
nariz arriba pujante le entró, a su padre así viendo.\\
Yendo de un salto hacia él, lo besó y lo abrazaba diciendo:\\
«Aquel por el que preguntas soy yo, padre mío, yo mismo,\\
que he regresado a la patria, después de veinte años ido.\\
Así que ten ya tu llanto, ten tú el lagrimoso gemido,\\
que ahora te tengo que hablar, y afanarse en el trance es preciso:\\
en nuestra casa he matado a los pretendientes altivos\\
dando a su ultraje doliente y a todos sus males castigo».
\end{versocita}

Laertes, claro, no se fía:

\begin{versocita}
Y le decía Laertes, alzando la voz, en respuesta:\\
«Si de verdad tú Odiseo, mi hijo tú, que aquí llegas,\\
dame ya seña bien clara para que yo me aconvenza».
\end{versocita}

Y Ulises le muestra la cicatriz que dejó el jabalí y que también ha reconocido antes su nodriza ---viene a ser, en el poema, como su DNI---.

\begin{versocita}
Y respondiendo decíale así el milardido Odiseo:\\
Pon esos ojos aquí, a esta herida no más lo primero,\\
la que me abrió con su blanco colmillo en Parnaso aquel puerco\\
cuando allí fui porque tú y la alta madre al abuelo materno,\\
a Autólico, me encomendasteis en pos de los dones aquellos\\
que prometió cuando vino y me confirmó con su gesto.
\end{versocita}

También le habla de los árboles del huerto, un pasaje harto más íntimo que el de la cansina cicatriz:

\begin{versocita}
Y, sí, también te diré de los árboles de este buen huerto,\\
los que me diste una vez, que uno a uno los iba pidiendo\\
yo cuando niño, y pegado a ti andaba detrás entre ellos\\
por el hortal; tú decías sus nombres y yo iba aprendiendo.\\
Trece perales me diste, y diez manzanos con ellos,\\
de higueras cuatro decenas, y liños en el viñedo\\
tú me darías cincuenta, en sazón cada cual a su tiempo,\\
que aquí racimos de toda sazón se levantan del suelo\\
cuando de arriba los cargan las estaciones de Zeus.
\end{versocita}

Y con esto, finalmente, padre e hijo se abrazan:

\begin{versocita}
[...] reconociendo señal tras señal que Odiseo le daba.\\
Echó a su hijo los brazos, y contra su cuerpo apretaba\\
el sufrimundo Odiseo a su padre, que allí desmayaba.
\end{versocita}

No recuerdo bien mi reacción a este episodio cuando lo leí por primera vez, probablemente mi estado emocional, después del asesinato de las niñas no daba mucho de sí. Pero al leerlo otra vez, a una edad a medio camino entre las que tendrían los dos hombres que se encuentran en el cuidado huerto de Laertes, no puedo evitar acordarme de mi padre y del padre de mi padre, a quien llamaban el tío Silvano, cabrero como Melantio ---perdóname entonces, compasiva lectora, caritativo lector, si me identifico más con el pobre desnarigado que con los niñatos que hemos dejado penando en el Hades---. De mi abuelo, recuerdo haberlo visto inclinado sobre su pequeño huerto en el pueblo, vestido, más o menos, como describe Homero, humilde, insignificante, callado. Recuerdo el cosquilleo en los labios cuando le besaba, por mucho que se rasurara, su barba era tan espesa y crecía tan deprisa que siempre había una fina capa, casi azul de tan oscura, poblando sus mejillas. Recuerdo, sobre todo, mi mano en la suya cuando me sacaba de paseo, algunas tardes. Era una mano callosa, de gruesos dedos, \emph{inmensa}, que contradecía su humildad y su silencio. Los callos los había labrado la azada, pero en mi imaginación infantil, no había diferencia con la espada de un general y ahora sé que no andaba desencaminado, ahora sé que Aquiles, el guerrero, envidiaba a mi abuelo, paseando a su nieto en un campo divino, bajo el bendito sol de primavera.

Mi padre, como mi abuelo, no era hombre de muchas palabras, ni de los que te abrazan cada vez que pueden ---mis pobres hijos, en cambio, tienen que soportar al pesado de su viejo achuchándoles con cualquier excusa--- ni era «mi mejor amigo», ni, que yo sepa, me dijo nunca «te quiero» con todas las letras ---también en eso somos muy diferentes, yo se lo repito a mis hijos hasta agotarlos---. Nunca intentó sobreprotegerme ni garantizarme un \emph{safe space} como el bueno de \emph{Nodiseo} a su hijito y sin embargo estaba siempre cerca. Quizás un poema me permita explicarme un poco. Se llama «Atajos»:

\begin{versopropio}
A mi padre le encantaban los atajos.\\
Intentaba descubrir uno nuevo\\
cada vez que venía a buscarme al colegio.
\stanzabreak
Normalmente, iba en el autobús\\
que unía Alba y Ebro,\\
diez kilómetros de carretera estrecha,\\
siempre atestada. A menudo pasábamos\\
una hora en el atasco, el conductor\\
maldiciendo en voz baja.
\stanzabreak
Se llamaba el tío Tralla.\\
Era grande y calvo. Yo siempre me sentaba\\
en la primera fila, a su lado.
\stanzabreak
Así había sido, desde la primera semana,\\
desde que Adrián y sus secuaces\\
descubrieron que acosar al chico extranjero\\
era un excelente pasatiempo. El tío Tralla\\
resolvió la situación de un solo manotazo,\\
luego señaló un asiento libre a su lado.
\stanzabreak
Estaba bien pensado. Yo era el primero en bajar,\\
y el tío Tralla se aseguraba\\
de que el grupo de Adrián fuese el último. Así que tenía\\
una buena ventaja hasta la seguridad del colegio.
\stanzabreak
La vuelta era más complicada,\\
había muchos sitios donde emboscarme,\\
camino de la parada del bus. A veces,\\
me libraba con solo unos cuantos pescozones,\\
pero de vez en cuando, Adrián tenía un mal día.
\stanzabreak
En casa,\\
mamá preguntaba por los moratones,\\
y yo fingía no oírla. Mamá se enfadaba mucho.\\
Decía que papá tenía que hacer algo al respecto.\\
Él preguntaba, ¿te duele, hijo?\\
Yo contestaba, No mucho.\\
Me daba una palmada en el hombro y miraba por la ventana.
\stanzabreak
Al día siguiente, venía a buscarme,\\
mi clase estaba en el primer piso,\\
y desde la ventana podía ver\\
su viejo Citroën aparcado frente al edificio.
\stanzabreak
A veces me esperaba mucho tiempo,\\
leyendo el libro que siempre llevaba en el bolsillo.
\stanzabreak
Cuando corría hacia él, siempre decía que era tarde.\\
Mamá nos estará esperando, aseguraba, tenemos que darnos prisa,\\
y mejor evitar la carretera congestionada.
\stanzabreak
Así que cada día intentaba un atajo diferente.
\stanzabreak
Es cierto, no había tráfico,\\
ni un alma en los caminos estrechos que elegía,\\
pero el atajo era tan enrevesado\\
que el coche no podía pasar entre los cañaverales. Entonces,\\
papá proponía que camináramos un poco.
\stanzabreak
Un día encontramos un zorro escabulléndose entre la hierba.\\
Otro, vimos peces voladores saltando en el río.
\stanzabreak
Nunca ahorrábamos un solo minuto.
\stanzabreak
Los atajos de mi padre eran tan largos como una tarde de verano,\\
esos largos veranos sin clases, sin bus y sin abusones.\\
Esos largos veranos en que volvíamos a nuestra ciudad,\\
la que habíamos tenido que dejar, la que una vez fue nuestra,\\
salvo que cada año lo era un poco menos.\\
Extraños en todas partes,\\
como siempre ocurre con los hijos de un marino.
\stanzabreak
Los atajos de mi padre eran tan complicados\\
que siempre nos perdíamos.\\
Entonces sacaba un bocadillo y una Coca-Cola,\\
y decía que un buen explorador siempre lleva comida y agua.
\stanzabreak
Tan perdido como decía estar,\\
siempre encontraba un claro soleado\\
para hacer nuestro picnic.
\stanzabreak
Tardábamos una eternidad en volver a casa,\\
pero mamá no parecía especialmente preocupada.\\
Yo lo sabía, porque la oía cantar,\\
y ella nunca cantaba si volvía en el autobús.
\stanzabreak
Aun así, miraba el reloj\\
y decía, ¿Cómo es que llegáis tan tarde?\\
Y yo decía, Papá intentó otro atajo,\\
y ella decía, A ver si algún día aprende.
\stanzabreak
Venía a recogerme,\\
incluso después de que Adrián se cansara del extranjero\\
y encontrara a otro de quién abusar.\\
Papá seguía apareciendo\\
con su viejo Citroën y su viejo libro.
\stanzabreak
Y decía que era tarde,\\
que mamá estaría esperando,\\
y que teníamos que darnos prisa,\\
pero no había problema, porque él conocía un atajo.
%\begin{versopropio}
%A papá le preocupaba que aprendiera a conducir,\\
%le preocupaba tanto que tomamos la costumbre\\
%de hacer prácticas mucho antes\\
%de tener la edad para sacarme el carnet.
%
%Aquellas escapadas,\\
%maravillosamente ilegales, mi yo adolescente,\\
%al timón de nuestro enorme coche familiar,\\
%más grande que el \emph{Nautilus}.
%
%Era también nuestro pequeño secreto,\\
%nos fugábamos discretamente,\\
%por carreteras soleadas que discurrían\\
%entre suburbios industriales y aldeas escondidas.
%
%Parábamos en pequeños cafés,\\
%donde charlábamos largamente. Después de todo,\\
%éramos socios al volante y lo suyo\\
%era intercambiar historias.
%
%Y mi padre tenía\\
%tantas historias que contarme.
%
%Cuando saqué el carnet siguió insistiendo\\
%en que practicar era imprescindible,\\
%declarándose todavía francamente inquieto\\
%con mi obvia falta de reflejos.
%
%Era una extraña obsesión. De hecho,\\
%yo había sacado el carnet a la primera\\
%y, después de tanta práctica, sabía muy bien\\
%que me sobraba habilidad. No importaba,
%
%papá declaró que todavía eran necesarias\\
%veinte mil leguas de viaje.
%
%Así que las hicimos. Viajábamos a veces\\
%al pueblo veraniego y pernoctábamos\\
%en la casita que teníamos allí\\
%antes de volver a casa al día siguiente.
%
%Papá cancelaba sus clases,\\
%anulaba compromisos, se saltaba citas,\\
%todo fuera por enseñarme\\
%a conducir, decía, como Dios manda.
%
%Y mamá, de repente,\\
%parecía igualmente preocupada\\
%y declaraba que aquellas excursiones\\
%eran absolutamente necesarias.
%
%Ayer pasé el día junto a mi hijo\\
%en otro largo viaje. Conduce tan bien\\
%como solía hacerlo yo, y aun así,\\
%prolongaría ese viaje hasta Finisterre.
%
%Viajaría con él veinte mil leguas, hasta llegar\\
%al final del Arco Iris, donde papá nos espera.
\end{versopropio}

Esa era mi Odisea diaria en la Ítaca que nunca fue, esos eran los trabajos de mi propio Laertes. Una palabra suya al director del instituto hubiera bastado para que a Adrián ---que no se llamaba Adrián igual que Alba no se llama Alba y el pueblo donde estudié el bachillerato elemental no se llama Ebro--- y sus abusones se les cayera el pelo, pero entonces su hijo no habría aprendido nunca a gestionar por su cuenta a los pretendientes.

Quizás por eso me ofenden tanto todas las astucias de Ulises con su padre. «¡Es tu viejo, infeliz!», me entran ganas de gritarle, «déjate de astucias y abrázale, ahora que puedes». Quizás por eso, el viejo Ulises, en su solitaria senectud, recuerda a su padre así:

\begin{versopropio}
Y Laertes más tarde. Te esperó,\\
y junto a ti celebró la carnicería de los jóvenes,\\
pero murió poco después. Ahora acude a tus sueños\\
a menudo, como si quisiera decirte algo esencial,\\
sobre la vida y la muerte. Algo que siempre,\\
simplemente escapa de tu comprensión.\\
Tú, que has matado a tantos hombres,\\
sin preguntar nunca por qué.
\stanzabreak
Sin embargo, tu padre era un hombre discreto,\\
y su fantasma no olvida sus modales.\\
No se demora demasiado, y su presencia no es fatigosa nunca.
\end{versopropio}

Alguna vez me han preguntado si me identifico con el Ulises de \emph{Abandonando Ítaca} y la respuesta es que él y yo tenemos en común todo el amor ---a Penélope, a Circe, a Telémaco--- del que el héroe fue capaz. Pero en otras cosas somos muy distintos. Y una de ellas tiene que ver con su Laertes y con el mío. El anciano Odiseo intenta entender qué quiere decirle su padre. Yo lo comprendí, para mi fortuna, hace muchos años.

De vuelta al canto XXIV, después del abrazo paterno-filial, los dos héroes pasan a asuntos prácticos, preocupándose por la posible revuelta de los familiares de los pretendientes masacrados. A partir de ahí se da por zanjada la parte sentimental y vuelven a aparecer los guerreros. En el tercer fragmento, Laertes toma parte en el último enfrentamiento entre Ulises y sus fieles y la banda de enojados padres a los que acaban de matarles a sus hijos. Y también en este caso me he encontrado yo siempre del lado equivocado. Si hubiera sido mi hijo el asesinado, me habría faltado tiempo para tomar una lanza y enfrentarme al asesino. Sin duda mi destino habría sido el mismo que el de Eupites, el padre de Antínoo, que se pone a la cabeza de la partida de rebeldes:

\begin{versocita}
Eupites los conducía igual que a niños pequeños;\\
y él, que la muerte del hijo pensaba vengar, nunca luego\\
ya tornaría, que allí lo aguardaba el destino funesto.
\end{versocita}

Me habría pasado como a Eupites, que muere al enfrentarse a Laertes, quien, por cierto, debe haber fumado algo muy fuerte: ha pasado de ser un viejo deshecho a estar en plena forma y sentir su pecho henchido de orgullo:

\begin{versocita}
[...] y en gozo Laertes a sí se decía:\\
«Dioses queridos, ¡qué día me dais! ¡Y cuánta alegría\\
ver a mi hijo y su hijo por la bravura en porfía!».
\end{versocita}

En realidad, no necesita fumar nada, tiene un camello divino para darle un chute de olímpica energía:

\begin{versocita}
Tal dijo; y Palas Atena alentole un vigor sin medida.\\
Y, haciendo rezo en voz alta a la de Zeus Cronida,\\
lanzó, blandiendo con pulso, su azcona larguisombría\\
y al casco le traspasó de Eupites la broncicarrilla\\
que de asta no lo libró, y entró el rejo de claro en su crisma:\\
él al caer retumbó, y sus armas chocáronle encima.
\end{versocita}

¡Pobre Eupites! Odiseo y los suyos caen sobre el resto de la banda, con la sana intención de cortarlos en cubitos, pero los olímpicos ya se han divertido bastante y decretan el fin de las hostilidades. Atenea da la voz de basta:

\begin{versocita}
«Teneos, hombres de Ítaca, de la agria guerra, y sin sangre\\
marchad cada uno a lo vuestro concordes y cuanto antes».
\end{versocita}

\begin{versocita}
Dijo Atenea, y a aquellos tomó la pavura amarilla;\\
y de sus manos, medrosos, volaban las armas, que iban\\
todas al suelo a parar, a la voz de la diosa que oían;\\
y se volvían a la ciudad anhelosos de vida.\\
El sufrimundo Odiseo rugió en malhorrísona grita\\
y prieto en sí se fue a ellos como águila de altanería.\\
Entonces suelta el Cronión una humosa centella encendida,\\
que vino a dar a los pies de la patrisoberbia ojiviva.
\end{versocita}

\begin{versocita}
Y dijo Atena ojiglauca, volviendo a Odiseo su vista:\\
«Sangre Laercia, Odiseo milmañas, hijo del cielo,\\
tente y pon fin al azote común de la guerra y su duelo,\\
para que no se te enoje Zeus Cronión miraluengo».
\end{versocita}

Así que la \emph{Odisea} acaba, literalmente, con un deus ex machina. Mi bienquerido director, dicho sea de paso, debería haberse estudiado el final del canto XXIV, con lo que nos habría ahorrado las ridículas frases grandilocuentes que pronuncian al alimón Penélope y Ulises mientras navegan hacia el sol poniente. Homero, en cambio, sabe que lo bueno, si breve. Después de la intervención de Atenea, cierra de un plumazo:

\begin{versocita}
Tal dijo Atena, y él bien la escuchó, en el alma alegrándose.\\
Y luego entre ellos dispuso los juros de paz perdurable\\
la diosa Palas Atena, la hija de Zeus cabrillante,\\
a Méntor siendo en figura y a par en voz semejante.
\end{versocita}

Y así acaba la \emph{Odisea}, la obra magna de Homero.
